\paragraph{C.3bb (Refuted): the divisor endpoint of integral finiteness.}
Let $k$ be any algebraically closed field of characteristic $p>0$.
For every $N\ge2$ there is an explicitly specified excellent normal
strictly henselian local $k$-algebra of dimension $N$, with residue
field $k$ and regular puncture $X$, such that
\[
 B^1(X)=\operatorname{Pic}(X)/
                  \operatorname{Pic}(X)_{\mathrm{alg}}
                          \longrightarrow(k,+)
\]
is surjective. The algebraic subgroup is defined by actual ordinary
curve families. Thus, together with C.3ba, integral finiteness fails
for every $N\ge2$ and every $1\le r\le N-1$ in every positive
characteristic. The exhibited quotient has exponent $p$; no rational
nonvanishing is asserted.

\begin{proof}
We supply the missing one-dimensional projective input explicitly;
the subsequent Segre argument does not require $r\ge2$.
Put $n=p(p+1)$ and take the plane curve, in coordinates $(Y:Z:W)$,
\[
 C:\quad Y^n-YW^{n-1}=Z^{n-1}W,\qquad
 v=(0:0:1),\qquad \infty=(0:1:0).
\]
On $W=1$ its equation is $y^n-y=z^{n-1}$ and its $y$
derivative is $-1$. On $W=0$ its only point is $\infty$, where
the $W$ derivative is $-1$. Therefore $C$ is smooth.
The standard plane-curve sequence gives $H^0(C,\mathcal O_C)=k$,
so it is connected and, being smooth, integral. Its genus is
\[
                         g=\frac{(n-1)(n-2)}2\ge10.
\]
Let $L=\Omega^1_{C/k}=\omega_C$. Its degree is
$2g-2=n(n-3)>0$, hence $L$ is ample. These are precisely
\citet[Lemmas 53.9.3, 53.8.4 and 33.44.14]{D_StacksCurveEndpoint}.
The identity of $L$ gives a canonical nowhere-zero section
\[
                \eta\in H^0(C,\Omega^1_{C/k}\otimes L^{-1}).
\]
On the affine chart, $z$ is an etale coordinate and $dz$ is a frame
of $L$. We use the corresponding fibre coordinate $t$ on $L^{-1}$;
thus the conormal frame $dt$ along the zero section corresponds to
$dz$. This specifies the normalization of $\eta$ without choosing
a global nowhere-zero regular differential on $C$.

Form the full section ring and the pointed strict henselization
\[
 A=\bigoplus_{q\ge0}H^0(C,L^q),\qquad
 R=(A_{A_+})^{\mathrm{sh}},\qquad
 Z=\operatorname{Tot}(L^{-1})\times_{\operatorname{Spec}A}
                                           \operatorname{Spec}R.
\]
The normality, finite generation, proper contraction, and isomorphism
away from the zero section follow from the full-section-ring argument
in C.3as, Step~3. In particular $Z$ is regular, the reduced exceptional
divisor is $i:C\hookrightarrow Z$, and $X=Z\setminus C$ is the
regular puncture of the normal excellent strict-local surface $R$.
No degree-one-generation assumption is needed.

The retraction $\pi:Z\to C$ splits the restricted differential sheaf:
\[
                   i^*\Omega^1_{Z/k}=\Omega^1_{C/k}\oplus L.
\]
Apply $d\log:\mathcal O_Z^*\to\Omega^1_{Z/k}$ to a line bundle,
restrict its cohomology class to $C$, and project to $L$. Serre trace
gives an integral group homomorphism
\[
 \rho:\operatorname{Pic}(Z)\longrightarrow H^1(C,L)
            =H^1(C,\omega_C)\xrightarrow{\operatorname{Tr}_{C/k}}k.
\]
The trace is dual to $H^0(C,\mathcal O_C)=k$ with its specified
unit, as in \citet[Lemma~48.27.1(5--7) and
Remarks~48.27.2--3]{D_Stacks0FVU}.
Tensor restriction is meant here: it retains the conormal generator
$dt$, rather than applying the differentials map that kills it.
Every line bundle pulled back from $C$ has zero conormal component.
In particular $\rho(\mathcal O_Z(C))=0$, since
$\mathcal O_Z(C)=\pi^*L^{-1}$. Localization therefore descends
$\rho$ to $\operatorname{Pic}(X)$.

We exhibit every scalar by actual divisors. For $b\in k$, impose
on the displayed affine total-space chart the equations
\[
                   z=bt,\qquad y^n-y=(bt)^{n-1}.
\]
The origin $y=t=0$ is smooth, since the derivative in $y$ is
$-1$, and hence lies on a unique irreducible component. Close that
component in the total space and normalize it. Its proper map to
its cone image is birational and has finite vertex fibre, hence is
finite. After pointed strict henselization, select the local factor
whose closed point is the specified germ $v$. The normalization is
an isomorphism at that smooth germ. Any remaining nonisomorphism
locus would be closed and proper over the local one-dimensional
branch, and so would meet its closed fibre. It is therefore empty.
This is the actual selected-branch argument of C.3as, Step~6, now
in dimension one. It gives a closed regular curve $D_b\subset Z$
meeting $C$ only at $v$, with local equation $z-bt$.

For the actual line bundle $\mathcal O_Z(D_b-D_0)$ the supported
$d\log$ class, restricted to $t=0$ while retaining $dt$, is
\[
 \left.\left(\frac{d(z-bt)}{z-bt}-\frac{dz}{z}\right)\right|_{t=0}
                          =-b\,\frac{dt}{z}.
\]
Under the chosen conormal frame this is the principal-part class
$-b\,dz/z$ in $H^1(C,\omega_C)$. Its trace is $-b$: the
local-to-global trace is dual to evaluation of constants at $v$,
so the coefficient is exactly one. There is no other contribution,
because the selected curves meet $C$ only at $v$. Replacing $\rho$
by its negative fixes the orientation and gives
\[
                      \rho(\mathcal O_Z(D_b-D_0))=b.
\]
The ordinary subgroup of $\operatorname{Pic}(X)$ is divisible by
$p$: extend each actual parameter-curve family to a smooth proper
curve by divisor closure, and use its action on the divisible group
$\operatorname{Pic}^0$ of that curve, exactly as in C.3as, Step~5.
Since $(k,+)$ has exponent $p$, $\rho$ kills that subgroup. This
proves the assertion for $N=2$.

For arbitrary $N\ge2$, put $m=N-2$ and use the specified data
\[
                 V_m=\mathbb P^m\times C,\qquad
                 L_m=\mathcal O(1)\boxtimes\omega_C,
\]
their full section ring, and its pointed strict henselization at the
vertex. Apply the construction and proof of C.3ba with $r=1$ and
the preceding $(C,L,D_b)$ as input. To make the endpoint explicit,
the map on resolutions is still multiplication by $z_0$, and the
arbitrary-coefficient free-basis and ind-etale factor proofs of
flatness are unchanged. The actual $D_b$ are Cartier divisors;
their flat inverse images remain Cartier divisors, regardless of
the fibre dimensions of an ind-etale localization.

The detector is now simply projection to the conormal cohomology
\[
 H^1(V_m,L_m)
       =H^0(\mathbb P^m,\mathcal O(1))\otimes H^1(C,\omega_C),
\]
followed by the trace on $C$ and extraction of the $z_0$ coefficient.
The exact coherent K\"unneth statement is
\citet[Lemma~33.29.1]{D_StacksSegreCohomology}.
The entire exceptional image is generated by the zero section and
has zero conormal component. On the chart $z_j\ne0$, the pulled-back
divisor equation is $z-b(z_0/z_j)t'_j$. The same computation gives
the scalar $b$ times the global section $z_0$, so the detector
attains $b$. Again it kills the $p$-divisible ordinary subgroup.
All pullbacks are performed on resolutions before localization;
no map between the two punctures is asserted.
This proves the result for every $N\ge2$.
\end{proof}

For $N=1$ the puncture has dimension zero and its positive-codimension
Chow groups vanish. The positive-characteristic integral finiteness
classification is therefore sharp in dimension and codimension.
The rational and ordinary-restriction questions remain as stated
elsewhere; this additive quotient settles neither of them.
